% =====================================================================
% IA en el Aula: guía práctica para profesores de la Universidad
% Panamericana. Taller interactivo (Beamer, 16:9).
% Compilar con pdflatex (dos pasadas). MiKTeX instala automáticamente
% fontawesome5, tcolorbox y listings si faltan.
% =====================================================================
\documentclass[aspectratio=169,10pt]{beamer}

\usepackage[T1]{fontenc}
\usepackage[utf8]{inputenc}
\usepackage[spanish,es-noshorthands]{babel}
\usepackage{tikz}
\usetikzlibrary{arrows.meta,calc,shadows,positioning}
\usepackage{amsmath,amssymb}
\usepackage{booktabs}
\usepackage{tabularx}
\usepackage{graphicx}
\usepackage[dvipsnames,svgnames,table]{xcolor}
\usepackage{fontawesome5}
\usepackage[most]{tcolorbox}
\tcbuselibrary{skins,breakable,listings}
\usepackage{listings}
\usepackage{hyperref}

% ---------------------------------------------------------------
% Paleta de colores
% ---------------------------------------------------------------
\definecolor{Dp}{RGB}{21,54,71}     % azul petróleo oscuro
\definecolor{Md}{RGB}{35,110,145}   % azul medio
\definecolor{Lt}{RGB}{120,185,205}  % azul claro
\definecolor{Ac}{RGB}{224,122,63}   % naranja (actividades)
\definecolor{Ac2}{RGB}{58,150,131}  % verde-azulado (tips)
\definecolor{Al}{RGB}{193,58,58}    % rojo (advertencias)
\definecolor{Gd}{RGB}{201,160,74}   % dorado
\definecolor{Bg}{RGB}{250,247,241}  % crema (fondo)
\definecolor{Wm}{RGB}{40,38,35}     % texto
\definecolor{Pr}{RGB}{110,80,150}   % morado (ámbito profesional)

% ---------------------------------------------------------------
% Tema
% ---------------------------------------------------------------
\usetheme{default}
\usecolortheme{default}
\usefonttheme{professionalfonts}
\setbeamertemplate{navigation symbols}{}
\setbeamercolor{normal text}{fg=Wm,bg=Bg}
\setbeamercolor{structure}{fg=Md}
\setbeamercolor{title}{fg=white,bg=Dp}
\setbeamercolor{frametitle}{fg=white,bg=Dp}
\setbeamercolor{section in toc}{fg=Dp}
\setbeamercolor{subsection in toc}{fg=Md}
\setbeamertemplate{section in toc shaded}[default][60]
\setbeamertemplate{subsection in toc shaded}[default][60]
\setbeamercolor{block title}{fg=white,bg=Md}
\setbeamercolor{block body}{fg=Wm,bg=Lt!20}
\setbeamercolor{itemize item}{fg=Ac}
\setbeamercolor{itemize subitem}{fg=Ac2}
\setbeamertemplate{itemize item}{\faCircle}
\setbeamertemplate{itemize subitem}{\faAngleRight}
\setbeamerfont{frametitle}{series=\bfseries,size=\Large}
\setbeamerfont{title}{series=\bfseries}
\setbeamertemplate{blocks}[rounded][shadow=false]
\setbeamertemplate{section in toc}[sections numbered]

\setbeamertemplate{footline}{%
  \leavevmode
  \hbox{\begin{beamercolorbox}[wd=\paperwidth,ht=2.4ex,dp=1.2ex,leftskip=1em,rightskip=1em]{normal text}
    {\color{Md!70}\scriptsize IA en el aula \textbullet\ Universidad Panamericana \hfill \insertframenumber/\inserttotalframenumber}
  \end{beamercolorbox}}%
}

% ---------------------------------------------------------------
% Cajas reutilizables
% ---------------------------------------------------------------
\newtcolorbox{actividadbox}[1][]{
  enhanced, breakable, colback=Ac!8, colframe=Ac, boxrule=1pt,
  arc=3mm, top=1.5mm, bottom=1.5mm, left=2mm, right=2mm,
  fonttitle=\bfseries\color{white}, coltitle=white, colbacktitle=Ac,
  title={\faUsers\ ACTIVIDAD \quad #1}
}
\newtcolorbox{tipbox}[1][EN RESUMEN]{
  enhanced, breakable, colback=Ac2!8, colframe=Ac2, boxrule=1pt,
  arc=3mm, top=1.5mm, bottom=1.5mm, left=2mm, right=2mm,
  fonttitle=\bfseries\color{white}, coltitle=white, colbacktitle=Ac2,
  title={\faLightbulb\ #1}
}
\newtcolorbox{warnbox}{
  enhanced, breakable, colback=Al!8, colframe=Al, boxrule=1pt,
  arc=3mm, top=1.5mm, bottom=1.5mm, left=2mm, right=2mm,
  fonttitle=\bfseries\color{white}, coltitle=white, colbacktitle=Al,
  title={\faExclamationTriangle\ CUIDADO}
}
\newtcolorbox{probox}{
  enhanced, breakable, colback=Pr!7, colframe=Pr, boxrule=1pt,
  arc=3mm, top=1.5mm, bottom=1.5mm, left=2mm, right=2mm,
  fonttitle=\bfseries\color{white}, coltitle=white, colbacktitle=Pr,
  title={\faBriefcase\ EN EL ÁMBITO PROFESIONAL}
}
\newcommand{\usocard}[3]{%
  \begin{tcolorbox}[enhanced,colback=white,colframe=Lt,boxrule=0.8pt,
    arc=2mm,left=2mm,right=2mm,top=1mm,bottom=1mm,
    title={\large #1\ \textbf{#2}},fonttitle=\bfseries\color{Dp},
    colbacktitle=Bg,coltitle=Dp]
    #3
  \end{tcolorbox}}
% Tarjeta de herramienta: icono, nombre, para qué sirve, cómo ayuda
\newcommand{\toolcard}[4]{%
  \begin{tcolorbox}[enhanced,colback=white,colframe=Md!50,boxrule=0.6pt,
    arc=2mm,left=1.5mm,right=1.5mm,top=0.8mm,bottom=0.8mm,
    title={#1\ \textbf{#2}},fonttitle=\small\bfseries\color{Dp},
    colbacktitle=Lt!25,coltitle=Dp]
    \scriptsize\textbf{\color{Md}Para qué sirve:} #3\\[1pt]
    \textbf{\color{Ac2}Cómo ayuda:} #4
  \end{tcolorbox}}

% Código LaTeX en diapositivas
\lstset{
  basicstyle=\ttfamily\scriptsize,
  keywordstyle=\color{Md}\bfseries,
  commentstyle=\color{Ac2}\itshape,
  stringstyle=\color{Ac},
  backgroundcolor=\color{white},
  frame=single, rulecolor=\color{Lt},
  breaklines=true, columns=fullflexible,
  showstringspaces=false,
  literate={á}{{\'a}}1 {é}{{\'e}}1 {í}{{\'i}}1 {ó}{{\'o}}1 {ú}{{\'u}}1 {ñ}{{\~n}}1 {¿}{{?`}}1
}

% Estilo de etapa con color como parámetro (definido fuera de los frames)
\tikzset{st/.style={draw=#1,fill=#1!10,rounded corners=2mm,minimum width=4.2cm,
  minimum height=2.3cm,align=left,text width=3.9cm,font=\scriptsize}}

\hypersetup{colorlinks=true, urlcolor=Ac, linkcolor=Md}

\title[IA en el aula]{Inteligencia Artificial en el Aula}
\subtitle{Guía práctica para profesores de la Universidad Panamericana}
\author[César Galindo]{César Galindo \\ \footnotesize Profesor Investigador}
\institute[UP]{Universidad Panamericana, Ciudad UP}
\date{\today}

\begin{document}

% ---------------------------------------------------------------
\begin{frame}[plain]
\begin{tikzpicture}[remember picture, overlay]
  \fill[Dp] (current page.south west) rectangle (current page.north east);
  \fill[Md!60] (current page.north east) circle (3.2);
  \fill[Ac!85] ([shift={(-0.6,0.6)}]current page.south west) circle (2.2);
  \fill[Ac2!70] ([shift={(1.4,-1.0)}]current page.north west) circle (1.1);
\end{tikzpicture}
\begin{center}
\vspace{0.7cm}
{\color{white}\Huge\bfseries Inteligencia Artificial\\[2pt] en el Aula}\\[10pt]
{\color{Gd!90}\Large Guía práctica para profesores}\\[2pt]
{\color{white!85}\large Claude, investigación profunda, Bloomberg, \LaTeX\ y GitHub}\\[18pt]
{\color{white}\faRobot\Large}\\[12pt]
{\color{white}\large César Galindo}\\
{\color{Gd!90}\normalsize Profesor Investigador}\\[2pt]
{\color{Lt}\normalsize Universidad Panamericana, Ciudad UP \quad\textbullet\quad \today}\\[4pt]
{\color{white!70}\small \href{mailto:jcgalindo@up.edu.mx}{jcgalindo@up.edu.mx}}
\end{center}
\end{frame}

% ---------------------------------------------------------------
\begin{frame}{Antes de empezar\dots}
\begin{center}
\Large\color{Dp}\faHandPaper
\end{center}
\vspace{8pt}
\begin{center}
\large\emph{``¿Cuántas veces han usado IA generativa este mes?''}
\end{center}
\vspace{10pt}
\begin{itemize}
  \item Levanten la mano: \textbf{nunca}.
  \item Levanten la mano: \textbf{1--2 veces}.
  \item Levanten la mano: \textbf{todos los días}.
  \item Y una más: \textbf{¿alguien ya usó Claude con su cuenta de la UP?}
\end{itemize}
\vspace{8pt}
\centering\small Cualquier respuesta es válida. Esta sesión es para decidir con información, no para convencer a nadie de nada.
\end{frame}

% ---------------------------------------------------------------
\begin{frame}{¿A dónde vamos hoy?}
\begin{columns}[T]
\column{0.55\textwidth}
\tableofcontents
\column{0.42\textwidth}
\begin{tipbox}[LO QUE SE LLEVAN]
\small
\begin{itemize}
  \item Un mapa de herramientas: qué hace cada una y cuándo conviene.
  \item Una receta para escribir buenas instrucciones.
  \item Un flujo de trabajo para exámenes en \LaTeX\ versionados en GitHub.
  \item Por qué todo esto importa para la vida profesional de sus alumnos.
\end{itemize}
\end{tipbox}
\end{columns}
\end{frame}

\AtBeginSection[]{
\begin{frame}{Agenda}
\tableofcontents[currentsection]
\end{frame}
}

% =====================================================================
\section{¿Qué es esto de la ``IA''?}
% =====================================================================

\begin{frame}{Una analogía, sin tecnicismos}
\begin{center}
\Large\color{Dp}\faComments
\end{center}
\vspace{6pt}
\begin{center}
\parbox{0.85\textwidth}{\centering\large
Un asistente como Claude, ChatGPT o Gemini es, en el fondo,\\[4pt]
\textbf{\color{Ac}un autocompletado extremadamente leído.}}
\end{center}
\vspace{10pt}
\begin{itemize}
  \item Leyó una fracción enorme de todo lo escrito en internet, libros y código.
  \item Aprendió qué palabra suele seguir a cuál, en qué contexto.
  \item Cuando le escriben algo, va prediciendo \emph{la continuación más probable}.
  \item Los asistentes actuales además \textbf{usan herramientas}: buscan en la web, leen PDFs, ejecutan cálculos, escriben código.
\end{itemize}
\vspace{6pt}
\begin{tipbox}
No ``entiende'' como una persona ni tiene criterio propio. Es estadística de lenguaje a una escala descomunal, más la capacidad de consultar fuentes.
\end{tipbox}
\end{frame}

% ---------------------------------------------------------------
\begin{frame}{Mitos vs.\ realidad}
\begin{columns}[T]
\column{0.5\textwidth}
\begin{tcolorbox}[enhanced,colback=Al!6,colframe=Al,arc=2mm,
  title={\faTimesCircle\ NO hace esto},fonttitle=\bfseries\color{white},colbacktitle=Al]
\begin{itemize}\small
  \item No sabe si algo es verdadero, solo si es \emph{plausible}
  \item No recuerda conversaciones pasadas, salvo que la herramienta lo active o se lo des tú
  \item No reemplaza tu criterio pedagógico ni disciplinar
  \item No garantiza que una cita o cifra exista
\end{itemize}
\end{tcolorbox}
\column{0.5\textwidth}
\begin{tcolorbox}[enhanced,colback=Ac2!6,colframe=Ac2,arc=2mm,
  title={\faCheckCircle\ SÍ hace esto},fonttitle=\bfseries\color{white},colbacktitle=Ac2]
\begin{itemize}\small
  \item Genera borradores rápidos de casi cualquier texto
  \item Adapta un contenido a distintos niveles
  \item Resume, reorganiza, traduce y reformula
  \item Escribe \LaTeX, código y fórmulas de hoja de cálculo
  \item Busca y resume fuentes (investigación profunda)
\end{itemize}
\end{tcolorbox}
\end{columns}
\vspace{8pt}
\begin{center}
\large\color{Dp}Trátenla como \textbf{borrador}, nunca como \textbf{producto final}.
\end{center}
\end{frame}

% ---------------------------------------------------------------
\begin{frame}{Las ``alucinaciones''}
\begin{warnbox}
La IA puede inventar datos, citas, fechas o nombres \textbf{con total seguridad}, sin avisar que no sabe. A esto se le llama \emph{alucinación}.
\end{warnbox}
\vspace{8pt}
\begin{center}
\begin{tcolorbox}[enhanced,colback=white,colframe=Dp!40,arc=2mm,width=0.9\textwidth]
\small\emph{``Cita tres estudios sobre aprendizaje basado en proyectos.''}\\[4pt]
$\rightarrow$ Puede devolver tres referencias que \textbf{suenan reales} pero no existen.\\[4pt]
$\rightarrow$ Las herramientas con búsqueda web reducen el riesgo porque \textbf{enlazan la fuente}, pero no lo eliminan: hay que abrir el enlace.
\end{tcolorbox}
\end{center}
\vspace{8pt}
\begin{center}
\large\color{Dp}\faSearch\ \color{Wm}\textbf{Todo lo que genera se revisa antes de usarse.}
\end{center}
\end{frame}

% =====================================================================
\section{El mapa de herramientas}
% =====================================================================

\begin{frame}{El mapa completo}
\begin{columns}[T]
\column{0.24\textwidth}
{\color{Md}\textbf{\faComments\ Asistentes generales}}
\begin{itemize}\small
  \item \href{https://claude.ai}{Claude} (Anthropic)
  \item \href{https://chatgpt.com}{ChatGPT} (OpenAI)
  \item \href{https://gemini.google.com}{Gemini} (Google)
  \item \href{https://copilot.microsoft.com}{Copilot} (Microsoft)
\end{itemize}
\column{0.24\textwidth}
{\color{Md}\textbf{\faChalkboardTeacher\ Docencia}}
\begin{itemize}\small
  \item \href{https://notebooklm.google}{NotebookLM}
  \item \href{https://www.magicschool.ai}{MagicSchool}
  \item \href{https://diffit.me}{Diffit}
  \item \href{https://curipod.com}{Curipod}
  \item \href{https://www.canva.com/magic-studio/}{Canva Magic Studio}
\end{itemize}
\column{0.24\textwidth}
{\color{Md}\textbf{\faSearchPlus\ Investigación}}
\begin{itemize}\small
  \item Research / Deep Research (Claude, ChatGPT, Gemini)
  \item \href{https://www.perplexity.ai}{Perplexity}
  \item \href{https://elicit.com}{Elicit}
  \item \href{https://consensus.app}{Consensus}
  \item \href{https://www.semanticscholar.org}{Semantic Scholar}
\end{itemize}
\column{0.26\textwidth}
{\color{Pr}\textbf{\faBriefcase\ Profesionales}}
\begin{itemize}\small
  \item Terminal \href{https://www.bloomberg.com/professional/products/bloomberg-terminal/}{Bloomberg} + ASKB
  \item \href{https://www.overleaf.com}{Overleaf} / \LaTeX
  \item \href{https://github.com}{GitHub} + Copilot
  \item Excel / Sheets con IA
\end{itemize}
\end{columns}
\vspace{8pt}
\begin{tipbox}
Las dos primeras columnas funcionan desde el navegador y no requieren programar. Las dos últimas son las que sus alumnos encontrarán en un despacho, un banco, una consultora o un laboratorio.
\end{tipbox}
\end{frame}

% ---------------------------------------------------------------
\begin{frame}{Asistentes generales: para qué sirve cada uno}
\begin{columns}[T]
\column{0.49\textwidth}
\toolcard{\faIcon{feather-alt}}{Claude}{redacción larga, análisis de documentos extensos (programas, reglamentos, artículos), \LaTeX\ y código.}{lee un PDF de 100 páginas y responde con citas al texto; genera exámenes directamente en \LaTeX.}
\toolcard{\faRobot}{ChatGPT}{uso general, imágenes, análisis de datos en hojas de cálculo.}{punto de partida conocido por casi todos los alumnos; útil para mostrar en clase qué ``ven'' ellos.}
\column{0.49\textwidth}
\toolcard{\faGoogle}{Gemini}{integración con Docs, Sheets, Gmail y Drive.}{si su flujo vive en Google Workspace, trabaja sobre sus archivos sin copiar y pegar.}
\toolcard{\faMicrosoft}{Copilot}{integración con Word, Excel, PowerPoint y Outlook.}{resume correos largos, propone fórmulas de Excel, arma una presentación a partir de un documento.}
\end{columns}
\end{frame}

% ---------------------------------------------------------------
\begin{frame}{Herramientas pensadas para docencia}
\begin{columns}[T]
\column{0.49\textwidth}
\toolcard{\faBookOpen}{NotebookLM}{responder \emph{solo} con base en los materiales que tú subes (apuntes, lecturas, programa).}{guías de estudio, preguntas frecuentes y resúmenes en audio de tu propio curso, con referencia a la página.}
\toolcard{\faChalkboardTeacher}{MagicSchool}{decenas de plantillas docentes: rúbricas, planeaciones, retroalimentación.}{resuelve tareas repetitivas con formularios; no hace falta escribir el prompt completo.}
\column{0.49\textwidth}
\toolcard{\faLayerGroup}{Diffit}{adaptar un mismo texto a distintos niveles de lectura.}{grupos heterogéneos, alumnos de intercambio, lecturas en inglés.}
\toolcard{\faPoll}{Curipod / Canva}{presentaciones interactivas, quizzes, organizadores gráficos.}{material visual en minutos que tú ajustas al grupo.}
\end{columns}
\end{frame}

% =====================================================================
\section{Claude en la UP: la cuenta institucional (plan Team)}
% =====================================================================

\begin{frame}{¿Qué cambia con una cuenta institucional?}
\begin{columns}[T]
\column{0.52\textwidth}
\begin{itemize}
  \item \textbf{Privacidad}: en los planes comerciales de Anthropic (Team, Enterprise, Education), las conversaciones \textbf{no se usan para entrenar modelos} por defecto.
  \item \textbf{Límites de uso más altos} que la cuenta gratuita: sesiones largas sin interrupciones.
  \item \textbf{Administración central}: la universidad gestiona altas, bajas y políticas.
  \item \textbf{Proyectos compartidos} entre colegas de una misma academia.
\end{itemize}
\column{0.45\textwidth}
\begin{tcolorbox}[enhanced,colback=white,colframe=Md,arc=2mm,
  title={\faUniversity\ Cuenta personal vs.\ UP},fonttitle=\bfseries\color{white},colbacktitle=Md]
\scriptsize
\begin{tabularx}{\linewidth}{@{}Xcc@{}}
\toprule
 & \textbf{Gratuita} & \textbf{Team UP}\\
\midrule
Entrena con tus datos & opcional & no\\
Límite de uso & bajo & alto\\
Proyectos compartidos & no & sí\\
Gestión institucional & no & sí\\
\bottomrule
\end{tabularx}
\end{tcolorbox}
\end{columns}
\vspace{6pt}
\begin{warnbox}
Aun con cuenta institucional: \textbf{no suban datos personales de alumnos} (nombres, matrículas, calificaciones). La regla de privacidad de la UP sigue mandando.
\end{warnbox}
\end{frame}

% ---------------------------------------------------------------
\begin{frame}{Las funciones de Claude que más sirven a un profesor}
\begin{columns}[T]
\column{0.49\textwidth}
\usocard{\faFolderOpen}{Proyectos}{\small Un espacio por materia: suben el programa, la bibliografía y sus apuntes, y escriben \emph{instrucciones fijas} (``responde como profesor de Finanzas Corporativas de 5º semestre, en español, con notación del libro de Berk--DeMarzo''). Cada conversación del proyecto parte de ese contexto.}
\usocard{\faShapes}{Artifacts}{\small Documentos, tablas, diagramas o pequeñas aplicaciones interactivas (simuladores, tarjetas de estudio) que aparecen a un lado del chat y se pueden compartir.}
\column{0.49\textwidth}
\usocard{\faSearchPlus}{Research}{\small Investigación profunda: Claude hace varias búsquedas, lee las fuentes y entrega un reporte con citas enlazadas.}
\usocard{\faPlug}{Conectores}{\small Acceso a Google Drive, Gmail, Calendar y otros servicios autorizados, para trabajar sobre sus propios archivos.}
\usocard{\faIcon{file-alt}}{Archivos}{\small Lee PDF, Word, Excel e imágenes; escribe \LaTeX, Markdown, Python o fórmulas.}
\end{columns}
\end{frame}

% ---------------------------------------------------------------
\begin{frame}{Ejemplo: un Proyecto por materia}
\begin{tikzpicture}[
  box/.style={draw=Md,fill=white,rounded corners=2mm,minimum width=3.1cm,minimum height=1.25cm,align=center,font=\small},
  arr/.style={-{Stealth[length=2.5mm]},thick,Ac}]
\node[box,fill=Lt!25] (in) {\faFileUpload\\Programa, apuntes,\\exámenes anteriores};
\node[box,right=0.9cm of in,fill=Md!15] (proj) {\faFolderOpen\ \textbf{Proyecto}\\``Probabilidad 2026''\\+ instrucciones fijas};
\node[box,right=0.9cm of proj,yshift=1.0cm] (o1) {\faClipboardList\ Banco de\\reactivos en \LaTeX};
\node[box,right=0.9cm of proj,yshift=-1.0cm] (o2) {\faComment\ Guía de estudio\\y retroalimentación};
\draw[arr] (in) -- (proj);
\draw[arr] (proj.east) -- (o1.west);
\draw[arr] (proj.east) -- (o2.west);
\end{tikzpicture}
\vspace{8pt}
\begin{tcolorbox}[enhanced,colback=Ac!6,colframe=Ac,arc=2mm]
\small\textbf{Instrucción fija de ejemplo:} \emph{``Eres asistente del curso Probabilidad (licenciatura, 3er semestre). Usa la notación de los apuntes adjuntos. Cuando generes ejercicios, entrégalos en \LaTeX\ con la clase \texttt{exam}, con solución dentro de \texttt{\textbackslash begin\{solution\}}. Nunca inventes referencias; si no sabes algo, dilo.''}
\end{tcolorbox}
\end{frame}

% =====================================================================
\section{Hablarle a la IA: el \emph{prompt}}
% =====================================================================

\begin{frame}{La receta de un buen \emph{prompt}}
\begin{center}
\Large\color{Dp}\textbf{Prompt} = la instrucción que le escribes a la IA
\end{center}
\vspace{10pt}
\begin{columns}[T]
\column{0.2\textwidth}\centering
\faUserTie\Large\\[4pt]\textbf{Rol}\\{\scriptsize ¿quién debe ser?}
\column{0.2\textwidth}\centering
\faTasks\Large\\[4pt]\textbf{Tarea}\\{\scriptsize ¿qué necesitas?}
\column{0.2\textwidth}\centering
\faInfoCircle\Large\\[4pt]\textbf{Contexto}\\{\scriptsize ¿para quién, qué nivel?}
\column{0.2\textwidth}\centering
\faIcon{file-alt}\Large\\[4pt]\textbf{Formato}\\{\scriptsize ¿cómo lo quieres?}
\column{0.2\textwidth}\centering
\faCheckDouble\Large\\[4pt]\textbf{Criterio}\\{\scriptsize ¿cómo sabrás que está bien?}
\end{columns}
\vspace{10pt}
\begin{tcolorbox}[enhanced,colback=Ac!6,colframe=Ac,arc=2mm]
\small\textbf{Ejemplo:} \emph{``Eres} [rol] \textbf{profesor de Microeconomía.} \emph{Necesito} [tarea] \textbf{6 reactivos de opción múltiple sobre elasticidad,} \emph{para} [contexto] \textbf{alumnos de 2º semestre de Economía,} \emph{en} [formato] \textbf{\LaTeX\ con la clase exam y la solución de cada uno.} [criterio] \textbf{Cada distractor debe corresponder a un error típico; indica cuál.''}
\end{tcolorbox}
\end{frame}

% ---------------------------------------------------------------
\begin{frame}{Prompt pobre vs.\ prompt bueno}
\begin{columns}[T]
\column{0.48\textwidth}
\begin{tcolorbox}[enhanced,colback=Al!6,colframe=Al,arc=2mm,title={\faTimesCircle\ Pobre},fonttitle=\bfseries\color{white},colbacktitle=Al]
\centering\emph{``Hazme un examen de finanzas.''}
\end{tcolorbox}
\vspace{6pt}
\small Sin nivel, sin tema, sin formato $\rightarrow$ resultado genérico e inútil.
\column{0.48\textwidth}
\begin{tcolorbox}[enhanced,colback=Ac2!6,colframe=Ac2,arc=2mm,title={\faCheckCircle\ Bueno},fonttitle=\bfseries\color{white},colbacktitle=Ac2]
\centering\footnotesize\emph{``Genera un examen de 8 preguntas sobre valor presente y TIR para 4º semestre de Finanzas: 5 de cálculo y 3 de interpretación, dificultad media, en \LaTeX\ (clase exam) con puntos por pregunta y soluciones paso a paso. Usa tasas y flujos realistas para México.''}
\end{tcolorbox}
\vspace{6pt}
\small Específico $\rightarrow$ resultado directamente compilable (después de revisarlo).
\end{columns}
\vspace{8pt}
\begin{tipbox}
Iteren: \emph{``la pregunta 3 es demasiado fácil; súbele un nivel''}, \emph{``cambia los datos para una versión B''}. La primera respuesta es solo el primer borrador.
\end{tipbox}
\end{frame}

% ---------------------------------------------------------------
\begin{frame}{Actividad 1}
\begin{actividadbox}[\faClock\ 5 minutos, en parejas]
\begin{enumerate}
  \item Piensen en una tarea real que hagan seguido (planear una clase, hacer un examen, escribir retroalimentación).
  \item Escriban un prompt con las 5 partes: \textbf{Rol + Tarea + Contexto + Formato + Criterio}.
  \item Pruébenlo en Claude (o en la herramienta que prefieran) y compártanlo con la pareja de al lado.
\end{enumerate}
\end{actividadbox}
\vspace{10pt}
\centering\large\color{Dp} Al terminar, 2 o 3 parejas leen su prompt y muestran el resultado.
\end{frame}

% =====================================================================
\section{Investigación profunda (\emph{deep research})}
% =====================================================================

\begin{frame}{¿Qué es la ``investigación profunda''?}
\begin{columns}[T]
\column{0.5\textwidth}
Un chat normal responde con lo que ya ``sabe''. Un modo de investigación profunda:
\begin{enumerate}\small
  \item descompone la pregunta en sub-preguntas,
  \item hace decenas de búsquedas en la web,
  \item lee las páginas y documentos encontrados,
  \item entrega un reporte estructurado \textbf{con citas enlazadas}.
\end{enumerate}
Tarda de minutos a media hora; no es instantáneo.
\column{0.47\textwidth}
\begin{tcolorbox}[enhanced,colback=white,colframe=Md,arc=2mm,
  title={\faSearchPlus\ Dónde encontrarla},fonttitle=\bfseries\color{white},colbacktitle=Md]
\footnotesize
\begin{itemize}
  \item \textbf{Claude}: botón \emph{Research}
  \item \textbf{ChatGPT}: \emph{Deep research}
  \item \textbf{Gemini}: \emph{Deep Research}
  \item \textbf{Perplexity}: respuestas rápidas con fuentes
  \item \textbf{Elicit}, \textbf{Consensus}, \textbf{Semantic Scholar}: literatura académica revisada por pares
\end{itemize}
\end{tcolorbox}
\end{columns}
\vspace{6pt}
\begin{tipbox}
Las herramientas generales buscan en la web abierta; Elicit, Consensus y Semantic Scholar buscan en bases de artículos científicos. Para una revisión de literatura, combinen ambas.
\end{tipbox}
\end{frame}

% ---------------------------------------------------------------
\begin{frame}{Un flujo de trabajo en tres etapas}
\begin{tikzpicture}[arr/.style={-{Stealth[length=2.5mm]},very thick,Dp}]
\node[st=Md] (a) {\textbf{\small 1. Descubrir}\\[2pt]Research / Deep Research o Perplexity para el panorama; Elicit o Consensus para artículos.\\[2pt]\emph{Salida:} lista de fuentes.};
\node[st=Al,right=0.6cm of a] (b) {\textbf{\small 2. Verificar}\\[2pt]Abrir cada enlace. ¿Existe? ¿Dice lo que el reporte afirma? ¿Es fuente primaria?\\[2pt]\emph{Salida:} fuentes confirmadas.};
\node[st=Ac2,right=0.6cm of b] (c) {\textbf{\small 3. Sintetizar}\\[2pt]Subir las fuentes confirmadas a un Proyecto de Claude o a NotebookLM y trabajar solo sobre ellas.\\[2pt]\emph{Salida:} lectura, caso o examen.};
\draw[arr] (a) -- (b);
\draw[arr] (b) -- (c);
\end{tikzpicture}
\vspace{10pt}
\begin{columns}[T]
\column{0.48\textwidth}
\usocard{\faChalkboardTeacher}{En docencia}{\small Actualizar un caso de estudio con datos de este año; preparar una lectura sobre una reforma reciente.}
\column{0.48\textwidth}
\usocard{\faFlask}{En investigación}{\small Mapear la literatura de un tema nuevo antes de leer a fondo; detectar artículos que citan a uno clave.}
\end{columns}
\end{frame}

% ---------------------------------------------------------------
\begin{frame}{Actividad 2}
\begin{actividadbox}[\faClock\ 8 minutos, individual]
\begin{enumerate}
  \item Elijan un tema de su materia que haya cambiado en los últimos dos años.
  \item Lancen una investigación profunda: \emph{``¿Qué ha cambiado en [tema] desde 2024? Dame un reporte de una página con fuentes primarias.''}
  \item Abran \textbf{tres} de las fuentes citadas. ¿Existen? ¿Dicen lo que el reporte afirma?
\end{enumerate}
\end{actividadbox}
\vspace{8pt}
\centering\large\color{Dp} Comentamos: ¿cuántas de las tres citas resistieron la verificación?
\end{frame}

% =====================================================================
\section{Bloomberg: datos financieros con IA}
% =====================================================================

\begin{frame}{La Terminal Bloomberg: qué es y para qué sirve}
\begin{columns}[T]
\column{0.52\textwidth}
\begin{itemize}\small
  \item La plataforma de datos de referencia en bancos, casas de bolsa, fondos y tesorerías: precios en tiempo real, estados financieros, bonos, divisas, materias primas, noticias y análisis.
  \item Se usa con \textbf{funciones} (``mnemónicos''): \texttt{DES} describe una empresa, \texttt{FA} sus estados financieros, \texttt{GP} grafica precios, \texttt{ECO} el calendario económico.
  \item Tiene complemento para \textbf{Excel} y \textbf{BQL / BQuant} para análisis en Python.
\end{itemize}
\column{0.45\textwidth}
\begin{tcolorbox}[enhanced,colback=white,colframe=Pr,arc=2mm,
  title={\faRobot\ ASKB: la IA dentro de Bloomberg},fonttitle=\bfseries\color{white},colbacktitle=Pr]
\footnotesize
\begin{itemize}
  \item Interfaz conversacional lanzada en 2026: preguntas en lenguaje natural sobre los datos y noticias de Bloomberg.
  \item Responde con \textbf{atribución a la fuente} dentro de la Terminal.
  \item Cuando hay análisis de datos, entrega el código BQL para llevarlo a Excel o BQuant.
\end{itemize}
\end{tcolorbox}
\end{columns}
\vspace{2pt}
\begin{probox}
\small Un analista le pregunta a ASKB, verifica la respuesta contra los datos y la lleva a un modelo. Ese criterio de verificación es lo que el aula debe entrenar.
\end{probox}
\end{frame}

% ---------------------------------------------------------------
\begin{frame}{Bloomberg en el aula: cómo ayuda}
\begin{columns}[T]
\column{0.49\textwidth}
\toolcard{\faCertificate}{Bloomberg Market Concepts (BMC)}{curso en línea de unas 8--12 horas sobre economía, divisas, renta fija y acciones, con certificado.}{tarea extracurricular con valor curricular; el profesor puede crear un código de grupo y ver el avance de sus alumnos.}
\toolcard{\faChartLine}{Datos reales para clase}{series de tasas, tipo de cambio, curvas de bonos, múltiplos de empresas mexicanas.}{ejercicios con datos de esta semana, no del libro de 2015.}
\column{0.49\textwidth}
\toolcard{\faFileExcel}{Excel + Bloomberg}{descargar datos a una hoja con fórmulas del complemento.}{los alumnos construyen modelos de valuación como en un banco.}
\toolcard{\faComments}{ASKB + verificación}{preguntar en lenguaje natural y contrastar con la función original.}{ejercicio de pensamiento crítico: ¿coincide la respuesta de la IA con \texttt{FA}?}
\end{columns}
\vspace{4pt}
\begin{tipbox}
Pregunten en su facultad por el acceso a terminales en los laboratorios de la UP. Sin terminal, BMC y los datos públicos de Banxico, INEGI y la BMV permiten ejercicios equivalentes.
\end{tipbox}
\end{frame}

% =====================================================================
\section{\LaTeX\ para preparar exámenes}
% =====================================================================

\begin{frame}{¿Por qué \LaTeX\ para exámenes?}
\begin{columns}[T]
\column{0.5\textwidth}
\begin{itemize}
  \item \textbf{Fórmulas impecables}: matemáticas, estadística, finanzas, física, química.
  \item \textbf{Versiones A/B} con el mismo archivo: cambian los datos, no el formato.
  \item \textbf{Soluciones en el mismo archivo}: una línea decide si se imprimen.
  \item \textbf{Puntos y tabla de calificación} automáticos.
  \item Texto plano: se versiona en GitHub y la IA lo escribe muy bien.
\end{itemize}
\column{0.47\textwidth}
\toolcard{\faLeaf}{Overleaf}{editor de \LaTeX\ en el navegador, sin instalar nada; colaboración en tiempo real.}{arrancar hoy mismo; compartir un examen con el coordinador de academia.}
\toolcard{\faIcon{feather-alt}}{Claude}{escribe y corrige \LaTeX: ``conviértelo a la clase exam'', ``¿por qué no compila?''.}{elimina la curva de aprendizaje inicial.}
\end{columns}
\end{frame}

% ---------------------------------------------------------------
\begin{frame}[fragile]{Un examen mínimo con la clase \texttt{exam}}
\begin{columns}[T]
\column{0.56\textwidth}
\begin{lstlisting}[language={[LaTeX]TeX}]
\documentclass[12pt]{exam}
\usepackage[spanish]{babel}
\usepackage{amsmath}
% \printanswers  % <- quita el % para la clave
\begin{document}
\begin{center}\textbf{Finanzas I -- Examen 1}\end{center}
Nombre: \underline{\hspace{6cm}}
\begin{center}\gradetable[h][questions]\end{center}
\begin{questions}
\question[4] Calcula el valor presente de
  \$10{,}000 a 3 años con $r=10\%$.
  \begin{solution}
    $VP = 10000/(1.1)^3 = 7513.15$
  \end{solution}
\question[2] La TIR es la tasa que:
  \begin{choices}
    \choice maximiza el VPN
    \CorrectChoice hace el VPN igual a cero
    \choice iguala la tasa libre de riesgo
  \end{choices}
\end{questions}
\end{document}
\end{lstlisting}
\column{0.41\textwidth}
\begin{tcolorbox}[enhanced,colback=white,colframe=Lt,arc=2mm,
  title={\faListOl\ Lo que hace cada parte},fonttitle=\bfseries\color{Dp},colbacktitle=Bg]
\scriptsize
\begin{itemize}
  \item \texttt{\textbackslash question[4]}: pregunta de 4 puntos.
  \item \texttt{\textbackslash gradetable}: tabla de calificación automática.
  \item \texttt{solution}: solo se imprime con \texttt{\textbackslash printanswers}.
  \item \texttt{\textbackslash CorrectChoice}: marca la correcta en la clave.
\end{itemize}
\end{tcolorbox}
\vspace{4pt}
\begin{tipbox}
\scriptsize Un archivo, dos PDF: el del alumno y la clave del profesor.
\end{tipbox}
\end{columns}
\end{frame}

% ---------------------------------------------------------------
\begin{frame}{Flujo de trabajo: de la idea al examen impreso}
\begin{tikzpicture}[
  st/.style={draw=Md,fill=white,rounded corners=2mm,minimum width=2.45cm,minimum height=1.6cm,align=center,text width=2.3cm,font=\scriptsize},
  arr/.style={-{Stealth[length=2mm]},thick,Ac}]
\node[st] (s1) {\faLightbulb\\[2pt]\textbf{Temario}\\objetivos y nivel};
\node[st,right=0.35cm of s1] (s2) {\faIcon{feather-alt}\\[2pt]\textbf{Claude}\\borrador en \LaTeX\ (clase exam)};
\node[st,right=0.35cm of s2,draw=Al] (s3) {\faSearch\\[2pt]\textbf{Revisión}\\resolver cada pregunta uno mismo};
\node[st,right=0.35cm of s3] (s4) {\faLeaf\\[2pt]\textbf{Overleaf}\\compilar versión A y B};
\node[st,right=0.35cm of s4] (s5) {\faGithub\\[2pt]\textbf{GitHub}\\guardar en repo privado};
\draw[arr] (s1) -- (s2); \draw[arr] (s2) -- (s3); \draw[arr] (s3) -- (s4); \draw[arr] (s4) -- (s5);
\end{tikzpicture}
\vspace{10pt}
\begin{columns}[T]
\column{0.48\textwidth}
\begin{tcolorbox}[enhanced,colback=Ac!6,colframe=Ac,arc=2mm]
\scriptsize\textbf{Prompt para versión B:} \emph{``Con este examen en \LaTeX, genera una versión B: mismos temas y dificultad, datos numéricos distintos, opciones en otro orden. Recalcula todas las soluciones y muéstrame el cálculo.''}
\end{tcolorbox}
\column{0.48\textwidth}
\begin{warnbox}
\scriptsize La IA se equivoca en aritmética y en detalles conceptuales. \textbf{Resuelvan cada pregunta ustedes} antes de aplicarla; la clave es responsabilidad del profesor.
\end{warnbox}
\end{columns}
\end{frame}

% =====================================================================
\section{GitHub para profesores}
% =====================================================================

\begin{frame}{GitHub: qué es y para qué sirve}
\begin{columns}[T]
\column{0.5\textwidth}
\begin{itemize}
  \item Plataforma para guardar archivos con \textbf{historial completo de cambios}: quién cambió qué, cuándo y por qué.
  \item Es el estándar de la industria para software, ciencia de datos y, cada vez más, documentación técnica.
  \item Repositorios \textbf{privados} (exámenes, claves) y \textbf{públicos} (apuntes, ejercicios, sitio del curso).
\end{itemize}
\column{0.47\textwidth}
\begin{tcolorbox}[enhanced,colback=white,colframe=Dp,arc=2mm,
  title={\faGraduationCap\ GitHub Education (gratis)},fonttitle=\bfseries\color{white},colbacktitle=Dp]
\footnotesize
Con correo \texttt{@up.edu.mx} y verificación como docente:
\begin{itemize}
  \item \textbf{GitHub Team} gratuito: repositorios privados ilimitados.
  \item \textbf{GitHub Copilot Pro} gratuito: asistente de IA para código y \LaTeX.
  \item \textbf{GitHub Classroom}: tareas por alumno, entregas y calificación automática.
\end{itemize}
\end{tcolorbox}
\end{columns}
\end{frame}

% ---------------------------------------------------------------
\begin{frame}{GitHub en el aula: cómo ayuda}
\begin{columns}[T]
\column{0.49\textwidth}
\toolcard{\faHistory}{Banco de exámenes versionado}{un repositorio privado por materia con todos los exámenes en \LaTeX.}{recuperar cualquier versión anterior; reutilizar y mejorar reactivos semestre a semestre.}
\toolcard{\faUsers}{Trabajo de academia}{varios profesores editan el mismo examen departamental.}{cada cambio queda registrado y se puede revisar antes de aceptarlo.}
\column{0.49\textwidth}
\toolcard{\faGlobe}{GitHub Pages}{publicar gratis el sitio de un curso desde un repositorio.}{programa, apuntes y ejercicios resueltos en una URL estable.}
\toolcard{\faTasks}{GitHub Classroom}{repartir y recoger tareas de programación o análisis de datos.}{los alumnos practican el flujo que usarán en su primer empleo.}
\end{columns}
\vspace{4pt}
\begin{warnbox}
\scriptsize Exámenes y claves \textbf{solo en repositorios privados}. Un repositorio público es visible para cualquiera, incluidos sus alumnos.
\end{warnbox}
\end{frame}

% =====================================================================
\section{Por qué importa: el ámbito profesional}
% =====================================================================

\begin{frame}{Lo que dicen los empleadores}
\begin{columns}[c]
\column{0.32\textwidth}\centering
{\Huge\color{Ac}\bfseries 39\%}\\[4pt]
\small de las habilidades centrales de los trabajadores cambiarán hacia 2030
\column{0.32\textwidth}\centering
{\Huge\color{Md}\bfseries 2/3}\\[4pt]
\small de los empleadores planean contratar talento con habilidades específicas de IA
\column{0.32\textwidth}\centering
{\Huge\color{Ac2}\bfseries 170\,M}\\[4pt]
\small empleos nuevos y 92\,M desplazados hacia 2030 (saldo neto: +78\,M)
\end{columns}
\vspace{12pt}
\begin{center}
\scriptsize Fuente: World Economic Forum, \emph{The Future of Jobs Report 2025}, encuesta a más de 1\,000 empleadores en 55 economías.
\end{center}
\vspace{6pt}
\begin{probox}
La pregunta para el aula ya no es si los egresados usarán IA, sino si sabrán \textbf{usarla con criterio}: verificar, citar, decidir cuándo no usarla.
\end{probox}
\end{frame}

% ---------------------------------------------------------------
\begin{frame}{Cada profesión, su herramienta}
\centering
\renewcommand{\arraystretch}{1.25}
\footnotesize
\begin{tabularx}{\textwidth}{@{}l>{\raggedright}X>{\raggedright\arraybackslash}X@{}}
\toprule
\textbf{Ámbito} & \textbf{Herramientas en uso} & \textbf{Qué entrenar en el aula}\\
\midrule
\faChartLine\ Finanzas y banca & Terminal Bloomberg con ASKB; Excel con IA & Contrastar la respuesta de la IA con el dato fuente\\
\faCode\ Ingeniería y datos & GitHub, Copilot, Claude, notebooks & Leer y probar código generado; control de versiones\\
\faBalanceScale\ Derecho & Asistentes legales con IA; investigación con fuentes & Verificar cada cita de jurisprudencia y ley\\
\faBriefcase\ Consultoría y empresa & Asistentes internos, investigación profunda & Distinguir dato, supuesto y opinión en un reporte\\
\faFlask\ Ciencia e investigación & \LaTeX, Overleaf, Elicit, Semantic Scholar & Revisión de literatura honesta; reproducibilidad\\
\faHeartbeat\ Salud & Resúmenes de evidencia clínica & Fuente primaria antes que resumen\\
\bottomrule
\end{tabularx}
\vspace{8pt}
\begin{tipbox}
En todos los renglones aparece lo mismo: \textbf{verificar}. Es la competencia transversal que más vale enseñar.
\end{tipbox}
\end{frame}

% ---------------------------------------------------------------
\begin{frame}{Actividad 3}
\begin{actividadbox}[\faClock\ 7 minutos, por academia]
\begin{enumerate}
  \item Elijan \textbf{una} herramienta de esta sesión que se use en la profesión de sus alumnos.
  \item Diseñen en 3 líneas una actividad de clase donde el alumno la use \emph{y además verifique} lo que produce.
  \item Decidan: ¿cómo se evalúa? ¿qué tiene que entregar el alumno además del resultado?
\end{enumerate}
\end{actividadbox}
\vspace{8pt}
\centering\large\color{Dp} Cada academia comparte su actividad en un minuto.
\end{frame}

% =====================================================================
\section{Riesgos, ética y límites}
% =====================================================================

\begin{frame}{Cinco riesgos que conviene conocer}
\begin{enumerate}
  \item \textbf{Alucinaciones} --- inventa datos y citas con total seguridad.
  \item \textbf{Privacidad} --- aun con cuenta institucional, nada de datos personales de alumnos.
  \item \textbf{Sesgos} --- reproduce los sesgos de los textos con los que fue entrenada.
  \item \textbf{Dependencia} --- usarla para \emph{todo} le quita al alumno (y al profesor) la oportunidad de pensar.
  \item \textbf{Confidencialidad} --- exámenes, claves y datos licenciados (como los de Bloomberg) no se publican ni se comparten fuera del uso autorizado.
\end{enumerate}
\vspace{8pt}
\begin{center}
\large\color{Dp}No usarla sigue siendo una opción válida. Si la usan, esto es lo que conviene vigilar.
\end{center}
\end{frame}

% ---------------------------------------------------------------
\begin{frame}{Honestidad académica}
\begin{center}
\Large\faGraduationCap\ \color{Dp}\textbf{¿Y los alumnos que la usan para hacer la tarea?}
\end{center}
\vspace{8pt}
\begin{itemize}
  \item Prohibirla por completo es una postura legítima; en la práctica es difícil de vigilar.
  \item Otra opción es \textbf{redefinir la tarea}: pedir proceso, no solo producto (borradores, fuentes verificadas, explicación oral, historial de cambios en GitHub).
  \item También se puede pedir que \textbf{declaren} si usaron IA, cuál y para qué.
  \item Los exámenes presenciales en papel, generados con \LaTeX, siguen siendo la evaluación más robusta.
\end{itemize}
\vspace{6pt}
\begin{tipbox}
La misma regla aplica al profesor: si usa IA para generar un examen o una rúbrica, la revisa línea por línea antes de aplicarla.
\end{tipbox}
\end{frame}

% =====================================================================
\section{Buenas prácticas y cierre}
% =====================================================================

\begin{frame}{Las 3 R}
\centering\small Una manera sencilla de usarla con cuidado:\\[10pt]
\begin{columns}[c]
\column{0.33\textwidth}\centering
{\Huge\color{Ac}\faSearch}\\[8pt]
{\Large\textbf{Revisar}}\\[4pt]
\small no usar un resultado sin leerlo por completo y abrir sus fuentes
\column{0.33\textwidth}\centering
{\Huge\color{Ac2}\faIcon{pencil-alt}}\\[8pt]
{\Large\textbf{Reescribir}}\\[4pt]
\small ajustar el tono, los datos y el nivel al grupo real
\column{0.33\textwidth}\centering
{\Huge\color{Md}\faBalanceScale}\\[8pt]
{\Large\textbf{Responsabilizarse}}\\[4pt]
\small lo que se entrega a los alumnos sigue siendo responsabilidad del profesor
\end{columns}
\end{frame}

% ---------------------------------------------------------------
\begin{frame}{El semáforo del uso responsable}
\begin{tcolorbox}[enhanced,colback=Ac2!8,colframe=Ac2,arc=2mm,title={\faCheckCircle\ Verde: sin problema},fonttitle=\bfseries\color{white},colbacktitle=Ac2]
\small Ideas, borradores, resúmenes, material de apoyo, \LaTeX\ para un examen que luego revisarán, mapas de literatura.
\end{tcolorbox}
\vspace{4pt}
\begin{tcolorbox}[enhanced,colback=Gd!15,colframe=Gd!70!black,arc=2mm,title={\faExclamationCircle\ Amarillo: revisar con cuidado},fonttitle=\bfseries\color{white},colbacktitle=Gd!70!black]
\small Datos, citas o cifras específicas; reportes de investigación profunda; respuestas de ASKB; claves de examen; cualquier cosa que se vaya a calificar.
\end{tcolorbox}
\vspace{4pt}
\begin{tcolorbox}[enhanced,colback=Al!8,colframe=Al,arc=2mm,title={\faTimesCircle\ Rojo: no hacerlo},fonttitle=\bfseries\color{white},colbacktitle=Al]
\small Subir nombres, matrículas, calificaciones o trabajos de alumnos reales; publicar exámenes o datos licenciados en repositorios públicos.
\end{tcolorbox}
\end{frame}

% ---------------------------------------------------------------
\begin{frame}{Recursos para seguir aprendiendo}
\begin{columns}[T]
\column{0.5\textwidth}
\begin{itemize}\small
  \item \faIcon{feather-alt}\ \href{https://claude.ai}{\textbf{Claude}} y su \href{https://anthropic.skilljar.com}{academia gratuita} de cursos
  \item \faBookOpen\ \href{https://notebooklm.google}{\textbf{NotebookLM}} --- tus materiales, tus respuestas
  \item \faSearchPlus\ \href{https://elicit.com}{\textbf{Elicit}}, \href{https://consensus.app}{\textbf{Consensus}}, \href{https://www.semanticscholar.org}{\textbf{Semantic Scholar}}
  \item \faChartLine\ \href{https://www.bloomberg.com/professional/products/bloomberg-terminal/education/certificate-courses/}{\textbf{Bloomberg Market Concepts}}
\end{itemize}
\column{0.5\textwidth}
\begin{itemize}\small
  \item \faLeaf\ \href{https://www.overleaf.com/learn}{\textbf{Overleaf Learn}} --- tutoriales de \LaTeX
  \item \faIcon{file-alt}\ \href{https://ctan.org/pkg/exam}{\textbf{Clase exam}} (CTAN, con manual)
  \item \faGithub\ \href{https://education.github.com/teachers}{\textbf{GitHub Education}} y \href{https://classroom.github.com}{\textbf{Classroom}}
  \item \faGlobe\ \href{https://www.weforum.org/publications/the-future-of-jobs-report-2025/}{\textbf{WEF Future of Jobs 2025}}
\end{itemize}
\end{columns}
\vspace{8pt}
\begin{tipbox}
Estas diapositivas, su código \LaTeX, una plantilla de examen y la guía de herramientas están en \href{https://github.com/IvliCaesar/ia-aula-profesores}{\textbf{github.com/IvliCaesar/ia-aula-profesores}}.
\end{tipbox}
\end{frame}

% ---------------------------------------------------------------
\begin{frame}{Cierre: ronda final}
\begin{center}
\Large\color{Dp}\faComments
\end{center}
\vspace{8pt}
{\large En una palabra, en voz alta:}
\begin{center}
\Large\color{Ac}\emph{``¿Qué herramienta van a probar esta semana, y para qué?''}
\end{center}
\vspace{8pt}
\centering\small Cada quien comparte con la persona de al lado; 2 o 3 lo dicen en voz alta para todo el grupo.
\end{frame}

% ---------------------------------------------------------------
\begin{frame}[plain]
\begin{tikzpicture}[remember picture, overlay]
  \fill[Dp] (current page.south west) rectangle (current page.north east);
  \fill[Ac!80] (current page.south east) circle (3.0);
  \fill[Ac2!60] ([shift={(0.8,0.8)}]current page.north west) circle (1.6);
\end{tikzpicture}
\begin{center}
\vspace{1.6cm}
{\color{white}\Huge\bfseries ¡Gracias!}\\[16pt]
{\color{Lt}\large César Galindo}\\
{\color{white!80}\normalsize Profesor Investigador \quad\textbullet\quad Universidad Panamericana, Ciudad UP}\\[2pt]
{\color{white!80}\normalsize \href{mailto:jcgalindo@up.edu.mx}{jcgalindo@up.edu.mx}}
\end{center}
\end{frame}

\end{document}
